\section*{第 \ref{ch:b3:hydrogeology} 章——水文地质学}

\begin{solution}{exo:b3:hydrogeology:1}
第一个滤水段：$z = 120 - 40 = \qty{80}{m}$，$h = 120 - 12.5 =
\qty{107.5}{m}$，压力\omterm{def:b3:hydrogeology:head}{水头} $107.5 - 80 = \qty{27.5}{m}$。第二个
滤水段：$z = \qty{40}{m}$，$h = \qty{105.0}{m}$，压力\omterm{def:b3:hydrogeology:head}{水头}
\qty{65.0}{m}。浅层滤水段的\omterm{def:b3:hydrogeology:head}{水头}较高，因此水
向下流动，\omterm{def:b3:hydrogeology:head}{水力梯度}为 $(107.5 - 105.0)/40 = \num{0.0625}$：
测压管位于补给区。仅凭压力
就会得出相反的结论。
\end{solution}

\begin{solution}{exo:b3:hydrogeology:2}
$q = K\Delta h/L = 12\times0.30/0.50 = \qty{7.2}{m/d}$。
截面积为 $\pi\times0.05^2 = \qty{7.85e-3}{m^2}$，故 $Q = \qty{0.0565}{m^3/d}$，即
$56.5/24 = \qty{2.36}{L/h}$。\omterm{def:b3:hydrogeology:seepage}{渗流速度}为 $7.2/0.35 =
\qty{20.6}{m/d}$。
\end{solution}

\begin{solution}{exo:b3:hydrogeology:3}
$K = 2.5/86400 = \qty{2.894e-5}{m/s}$，并由 $k = K\mu/(\rho g) =
\num{2.894e-5}\times\num{1.306e-3}/(999.7\times9.81) =
\qty{3.85e-12}{m^2}$。在 \qty{20}{\celsius} 时，$K$ 乘以
$(998.2/\num{1.002e-3})/(999.7/\num{1.306e-3}) = \num{1.30}$：$K =
\qty{3.25}{m/d}$。这一变化几乎全部源于动力黏度的变化。
\end{solution}

\begin{solution}{exo:b3:hydrogeology:4}
$T = 2.5\times180 = \qty{450}{m^2/d}$。水的膨胀贡献
$\num{9e-7}\times180 = \num{1.62e-4}$，约占 \qty{80}{\percent}，$S =
\num{2e-4}$；其余由石灰岩的压缩提供。
\end{solution}

\begin{solution}{exo:b3:hydrogeology:5}
由\cref{prop:b3:hydrogeology:thiem}，$T = Q\ln(r_2/r_1)/(2\pi(s_1 -
s_2)) = 1500\times\ln10/(2\pi\times1.53) = \qty{359}{m^2/d}$，于是 $K =
359/40 = \qty{9.0}{m/d}$。
\end{solution}

\begin{solution}{exo:b3:hydrogeology:6}
$Q/4\pi T = 1200/(4\pi\times300) = \qty{0.318}{m}$。经过 1 天，$u =
300^2\times\num{5e-4}/(4\times300\times1) = \num{0.0375}$，$W = -0.5772
+ 3.2834 + 0.0375 - 0.0004 = \num{2.743}$：$s = \qty{0.873}{m}$。经过
10 天后，$u = \num{0.00375}$，$W = -0.5772 + 5.5860 + 0.0038 = \num{5.013}$：
$s = \qty{1.596}{m}$。时间延长为 10 倍，
降深增加不到原来的两倍。
\end{solution}

\begin{solution}{exo:b3:hydrogeology:7}
$u \le 0.05$，即 $t \ge r^2S/(4T\times0.05)$。在 \qty{60}{m} 处：$3600\times
\num{2e-4}/90 = \qty{0.008}{d}$，\qty{11.5}{min}。在 \qty{500}{m} 处：
\qty{0.556}{d}，\qty{800}{min}，超过 13 小时。时间随 $r^2$ 增大：
在一两天的试验中，用于作直线的观测井
必须位于抽水井几百米之内。
\end{solution}

\begin{solution}{exo:b3:hydrogeology:8}
在中点处，两井各距 \qty{200}{m}：$u = 200^2\times\num{2e-4}/(4
\times450\times10) = \num{4.44e-4}$，$W = \num{7.142}$，$s = 0.382\times
7.142 = \qty{2.728}{m}$（每口井），合计 \qty{5.456}{m}。在
\qty{400}{m} 处，$u = \num{1.78e-3}$，$W = \num{5.757}$，$s =
\qty{2.199}{m}$。每口井使另一口井的水位下降
\qty{2.2}{m}：其水泵须将水提升 \qty{2.2}{m}，这
需要耗费能量，并且在泵吸水口干涸之前，剩余的
可用降深也更小。
\end{solution}

\begin{solution}{exo:b3:hydrogeology:9}
坝底与基底之间的四条流线围成 $N_f = 5$ 个
流道；九条等势线将 \qty{10}{m} \omterm{def:b3:hydrogeology:head}{水头}分为 $N_d = 10$ 个
等势降。$q = KHN_f/N_d = 5\times10\times5/10 = \qty{25}{m^3/d}$，即每米坝长
的流量；坝长 \qty{300}{m} 时为 \qty{7500}{m^3/d}。计算得到的
形状因子为 0.512 而非 0.5，则每米坝长给出 \qty{25.6}{m^3/d}。
\end{solution}

\begin{solution}{exo:b3:hydrogeology:10}
$k = na^2/8 = 0.30\times(\num{5e-5})^2/8 = \qty{9.4e-11}{m^2}$，并由 $K =
k\rho g/\mu = \num{9.375e-11}\times998.2\times9.81/\num{1.002e-3} =
\qty{9.2e-4}{m/s}$，即 \qty{79}{m/d}，处于干净砂的范围 0.3 至
\qty{150}{m/d} 之间。模型忽略了路径的迂曲度，路径比直管长（使 $K$ 降低）、
孔隙之间的喉道，比孔隙窄（同样使其降低）、
两端封闭、不参与流动的盲孔（同样使其降低）以及
孔径分布的分散性：由于管道的流量与 $a^4$ 成正比，
最宽的路径输送了大部分水量，
孔径分布的分散使 $K$ 高于平均孔径管道的值。
\end{solution}

\begin{solution}{exo:b3:hydrogeology:11}
(a)~对数项在 $2.25\,Tt = r^2S$ 时为零，即 $r_0 =
\sqrt{2.25\,Tt/S}$。(b)~\omterm{def:b3:hydrogeology:aquifer-types}{承压含水层}与\omterm{def:b3:hydrogeology:aquifer-types}{潜水含水层}两种情形：
\[
r_0 = \sqrt{\frac{2.25\times450\times30}{\num{2e-4}}} = \qty{12.3}{km},
\qquad r_0 = \sqrt{\frac{2.25\times450\times30}{0.2}} = \qty{390}{m},
\]
小 $\sqrt{1000} = 31.6$ 倍。(c)~压力
以扩散系数 $T/S$ 扩散，$S$ 较小时扩散极快：在
试验期间，\omterm{def:b3:hydrogeology:cone}{降落漏斗}会波及断层、含水层的露头以及河流，
降深会偏离直线，这说明水文地质学家
已遇到了边界。
\end{solution}

\begin{solution}{exo:b3:hydrogeology:12}
(a)~在随水流运动的参考系中，前缘的标准差为
$\sigma_x = \sqrt{2Dt} = \sqrt{2\alpha_Lvt}$。前缘通过
$x$ 处约在 $t = x/v$ 时刻，此时 $\sigma_x = \sqrt{2\alpha_Lx}$。以
速度 $v$ 运动时，距离上的标准差 $\sigma_x$ 相当于
时间上的标准差 $\sigma_t = \sigma_x/v$。(b)~\qty{16}{\percent} 与 \qty{84}{\percent} 对应的时刻
位于均值两侧各一个标准差处：$\sigma_t = (247.2 -
152.6)/2 = \num{47.3}\text{ 天}$。于是 $\alpha_L = v^2\sigma_t^2/(2x) =
0.25\times47.3^2/200 = \qty{2.8}{m}$。（曲线是用
\qty{3}{m} 计算的：正态近似略微低估了该值。）(c)~$R =
1 + 1.9\times0.32/0.30 = \num{3.03}$：溶质运移速度
约为水流速度的三分之一。
\end{solution}

\begin{solution}{pb:b3:hydrogeology:1}
\textbf{1.} 页岩 $363.2 - 282.5 = \qty{80.7}{m}$；石灰岩 $282.5 -
102.5 = \qty{180}{m}$。

\textbf{2.} $h = 363.2 - 45.0 = \qty{318.2}{m}$；顶部的压力\omterm{def:b3:hydrogeology:head}{水头}为
$318.2 - 282.5 = \qty{35.7}{m}$。水位高出石灰岩
顶面 \qty{35.7}{m}，受页岩压制：该含水层为\omterm{def:b3:hydrogeology:aquifer-types}{承压含水层}。

\textbf{3.} 由 $h_2 - h_1 = 550a - 100b = -1.62$ 及 $h_3 - h_1 = 250a
- 600b = -1.44$：$a = \num{-2.715e-3}$ 与 $b = \num{1.269e-3}$，单位
为每米。

\textbf{4.} $|\nabla h| = \sqrt{a^2 + b^2} = \num{0.0030}$。水流沿
$-\nabla h = (\num{2.715e-3}, \num{-1.269e-3})$ 的方向流动，即向东、向南：
方位角 $\ang{180} - \arctan(2.715/1.269) = \ang{115}$。

\textbf{5.} 约 \qty{10}{min} 之前的数据点位于直线上方：此处
$u > 0.05$，泰斯曲线高于其对数近似
（$W > -\gamma - \ln u$，如井函数图所示）。

\textbf{6.} 自 \qty{15}{min} 起的数据点，最佳拟合直线经过
\qty{100}{min} 时的 \qty{1.75}{m} 与 \qty{1000}{min} 时的 \qty{2.63}{m}：
$\Delta s = \qty{0.87}{m}$ 每个对数循环。

\textbf{7.} $T = 2.3\times2160/(4\pi\times0.872) = \qty{453}{m^2/d}$。

\textbf{8.} 直线每循环下降 \qty{0.87}{m}，自 \qty{100}{min} 时的 \qty{1.75}{m} 起算，
零值在 $1.75/0.872 = 2.01$ 个循环之前即已到达，
$t_0 = \qty{0.98}{min} = \qty{6.83e-4}{d}$。

\textbf{9.} $S = 2.25\,Tt_0/r^2 = 2.25\times453\times\num{6.83e-4}/3600 =
\num{1.9e-4}$，处于\omterm{def:b3:hydrogeology:aquifer-types}{承压含水层}的数值范围内。

\textbf{10.} 在 \qty{15}{min} 时，$u = 3600\times\num{1.9e-4}/(4\times453
\times0.0104) = \num{0.036}$，之后更小。

\textbf{11.} $K = 453/180 = \qty{2.5}{m/d}$，为裂隙石灰岩。

\textbf{12.} $u = 3600\times\num{2e-4}/(4\times450\times2) =
\num{2.0e-4}$，$W = -0.5772 + 8.5172 + 0.0002 = \num{7.940}$，$s =
0.382\times7.940 = \qty{3.03}{m}$，与记录一致。

\textbf{13.} $u = 10^6\times\num{2e-4}/(4\times450\times30) =
\num{0.0037}$，$W = \num{5.02}$，$s = \qty{1.92}{m}$。

\textbf{14.} $r_0 = \sqrt{2.25\times450\times30/(\num{2e-4})} =
\qty{12.3}{km}$，超过了区域的宽度。\omterm{def:b3:hydrogeology:cone}{降落漏斗}已经到达
石灰岩的边缘、其露头与断层，断层或使其与渗透性较差的岩石相接，
因而泰斯的无限含水层模型已不再适用：
只有降深记录才能说明，
边界是补给水井还是阻断了水流。

\textbf{15.} $2160\times30 = \qty{64800}{m^3}$。该水量来自储存：
水体膨胀与石灰岩压缩释放的水，分布于整个
\omterm{def:b3:hydrogeology:cone}{降落漏斗}范围。在\omterm{def:b3:hydrogeology:cone}{降落漏斗}上对 $S\,s$ 积分，
即得 $Qt$。

\textbf{16.} $u = 10^6\times0.18/(4\times450\times30) = \num{3.33}$，$W
= \num{0.0086}$，$s = \qty{3}{mm}$：农场不会察觉。
抽水将取用井附近孔隙中的水，而不是
整个含水层的压力。

\textbf{17.} $q = 2.5\times0.0030 = \qty{0.0075}{m/d}$；$v = 0.0075/0.18
= \qty{0.042}{m/d}$，即每年 \qty{15}{m}。

\textbf{18.} $200/0.0417 = \num{4800}\text{ 天}$，即 \num{13.1}~年；污染物
则需 \num{32.9}~年。

\textbf{19.} 通过半径为 $r$ 的圆柱面，流量 $Q$ 穿过
面积 $2\pi rb$，其中孔隙所占比例为 $n_e$：$\mathrm dr/\mathrm
dt = -Q/(2\pi rbn_e)$，将 $2\pi bn_er\,\mathrm dr = -Q\,
\mathrm dt$ 从 $r$ 积分到 0，得 $t = \pi n_ebr^2/Q$。此处 $\pi\times0.18
\times180\times200^2/2160 = \num{1885}\text{ 天}$，对水而言为 \num{5.2}\text{ 年}，
对污染物则为 $2.5\times5.2 = \num{12.9}\text{ 年}$：水井
使其提早了一倍以上。

\textbf{20.} $\alpha_L = \qty{20}{m}$，$\sqrt{2\times20\times200} =
\qty{89}{m}$。浓度达到其终值的 \qty{16}{\percent}，
即位于平均前缘前方 \qty{89}{m} 处：该浓度水平到达水井
时平均前缘仍距其 \qty{89}{m}，此时为 $(111/200)^2\times12.9 = \num{4.0}\text{ 年}$，
约为平均时间的三分之一。首批到达的痕迹关系到饮用水安全，
而非平均值。

\textbf{21.} 赞成：水井捕获了污染物，使其无法继续向下游扩散，
抽出的水可在抽水处进行处理。
反对：抽水使污染物比其自然扩散更快地进入供水井，
且供水须在数十年内
持续监测并处理。

\textbf{22.} 瓦尔登石灰岩的\omterm{def:b3:hydrogeology:storage}{导水系数}为 $T \approx
\textbf{\qty[text-series-to-math=true]{450}{m^2/d}}$（\qty{453}{m^2/d}
由试验得出）。
\end{solution}
